% ============================================================
% QUESTION TYPE: MULTIPLE CHOICE
%
% PEANUT EXAM MCQ LAYOUT STANDARD
%
% AI INSTRUCTION:
% - Follow this format EXACTLY for multiple-choice questions
% - Questions are grouped by PART using \section*
% - Individual questions are bold-numbered (not subsection)
% - Choices use enumerate with label=\alph*)
% - Use \vspace{...} for spacing between questions
% - Do NOT use \subsection or custom environments
% - Do NOT alter choice labels or spacing style
% ============================================================

\section*{Part I: [PART TITLE]}
\textit{Choose the best answer for each question.}

\vspace{5mm}

\textbf{1. This is a multiple-choice question statement.
It may span one or more lines and should be written clearly.}

\begin{enumerate}[label=\alph*)]
\item First option text
\item Second option text
\item Third option text
\item Fourth option text
\end{enumerate}

\vspace{15pt}

\textbf{2. This question may test definitions, concepts, or recognition.}

\begin{enumerate}[label=\alph*)]
\item Option A
\item Option B
\item Option C
\item Option D
\end{enumerate}

\vspace{15pt}

\textbf{3. This question may include inline math such as $x^2 + 1$ or
ask about structure, grammar, or interpretation.}

\begin{enumerate}[label=\alph*)]
\item Choice 1
\item Choice 2
\item Choice 3
\item Choice 4
\end{enumerate}

\vspace{15pt}

% ============================================================

\newpage
\textbf{4. This question continues numbering across pages.
New pages do NOT reset question numbers.}

\begin{enumerate}[label=\alph*)]
\item Option one
\item Option two
\item Option three
\item Option four
\end{enumerate}

\vspace{15pt}

\textbf{5. This question may check application or scenario-based reasoning.}

\begin{enumerate}[label=\alph*)]
\item Answer choice A
\item Answer choice B
\item Answer choice C
\item Answer choice D
\end{enumerate}

\vspace{15pt}
